\chapter{Réalisation}
\label{chap:realisation}

\section*{Introduction}
\addcontentsline{toc}{section}{Introduction}

Ce chapitre décrit la mise en œuvre du système, de l'extraction du corpus jusqu'aux interfaces.
Une place particulière y est accordée aux difficultés rencontrées : les défauts d'extraction
décrits en section~\ref{sec:bugs} ont représenté une part importante du temps de travail, et
illustrent une leçon qui vaut au-delà de ce projet --- la qualité du corpus conditionne tout ce
qui est construit dessus, et ses défauts ne produisent aucun message d'erreur.

\section{Constitution du corpus}

\subsection{Extraction et découpage}

Le texte a été extrait du PDF officiel du Code du travail à l'aide de PyMuPDF. Le découpage en
articles repose sur une expression régulière ancrée en début de ligne, appliquée au texte
\emph{complet} plutôt que page par page --- un article pouvant s'étendre sur deux pages, un
traitement page à page l'aurait systématiquement coupé.

Chaque article est enrichi de sa position exacte dans la hiérarchie de la loi (livre, titre,
chapitre, section), reconstruite en parcourant le texte dans l'ordre : les en-têtes de hiérarchie
annoncent ce qui suit, ils sont donc mémorisés au fil de la lecture et appliqués aux articles
suivants jusqu'au prochain en-tête.

Le corpus final compte \textbf{\NbArticles{} articles} répartis sur \NbLivres{} livres
(figure~\ref{fig:corpus}).

\begin{figure}[htbp]
  \centering
  % Structure du corpus. Le déséquilibre est celui de la loi elle-même : le
% Livre II concentre à lui seul près de la moitié des articles, ce qui
% explique que le jeu de référence y puise autant.
% Répartition des articles du corpus par Livre
% Généré par rapport/build_data.py — NE PAS ÉDITER.
\newcommand{\CorpusLabels}{{Livre II},{Livre I},{Livre III},{Livre IV},{Livre VI},{Livre V},{Livre Prélim.},{Livre VII}}
\newcommand{\CorpusCoords}{(261,0) (122,1) (79,2) (55,3) (37,4) (19,5) (12,6) (4,7)}
\newcommand{\CorpusMax}{261}

\begin{tikzpicture}

\begin{axis}[
  name=livres,
  barresH,
  height=5.4cm, width=0.55\linewidth,
  xmin=0, xmax=\CorpusMax*1.28,
  ytick={0,...,7},
  % Libellés en clair : pgfplots ne développe pas une macro donnée à
  % « yticklabels ». L'ordre suit celui des barres, par effectif décroissant.
  yticklabels={{Livre II},{Livre I},{Livre III},{Livre IV},
               {Livre VI},{Livre V},{Livre prélim.},{Livre VII}},
  yticklabel style={font=\sffamily\scriptsize, text=ink, align=right},
  xlabel={nombre d'articles},
  nodes near coords,
  bar width=9pt,
  title={Répartition par Livre},
  title style={at={(0,1.06)}, anchor=south west},
]
\addplot[draw=inptblue, fill=inptblue, fill opacity=0.85]
  coordinates {\CorpusCoords};
\end{axis}

% Décomposition du Livre II — le seul dont la taille justifie qu'on l'ouvre.
\begin{axis}[
  barresH,
  at={($(livres.north east)+(1.7cm,0)$)}, anchor=north west,
  height=5.4cm, width=0.32\linewidth,
  xmin=0, xmax=118,
  ytick={0,...,4},
  yticklabels={{Durée du travail},{Hygiène et sécurité},{Salaire},
               {Mineur et femme},{Dispositions générales}},
  yticklabel style={font=\sffamily\scriptsize, text=ink, align=right},
  xlabel={articles},
  nodes near coords,
  bar width=8pt,
  title={Livre II, par Titre},
  title style={at={(0,1.06)}, anchor=south west},
]
\addplot[draw=inptblue, fill=inptblue, fill opacity=0.45]
  coordinates {(97,0) (64,1) (51,2) (41,3) (8,4)};
\end{axis}

% Le commentaire sur les métadonnées est passé en légende : posé dans la
% figure, il débordait de la justification.

\end{tikzpicture}

  \caption[Structure du corpus extrait]{Structure du corpus extrait. Le déséquilibre est celui de
  la loi elle-même : le Livre~II concentre à lui seul 261 des \NbArticles{} articles, ce qui
  explique que le jeu de référence du chapitre~\ref{chap:evaluation} y puise une large part de ses
  questions. Chaque article porte en métadonnée son chemin complet --- livre, titre, chapitre,
  section --- reconstruit à la lecture ; ce chemin n'est pas vectorisé, il sert au filtrage et à
  l'affichage des sources, pas à la mesure de similarité.}
  \label{fig:corpus}
\end{figure}

Le résultat est sérialisé au format JSONL, un article par ligne :

\begin{lstlisting}[caption={Structure d'un article du corpus},label={lst:article}]
{
  "article_number": "152",
  "article_text": "La salariee en etat de grossesse attestee par certificat
                   medical dispose d'un conge de maternite de quatorze
                   semaines, sauf stipulations plus favorables...",
  "amende_2021": false,
  "livre": "Livre II: Des conditions de travail et de la remuneration du salarie",
  "titre": "Titre II : De la protection du mineur et de la femme",
  "chapitre": "Chapitre II : De la protection de la maternite",
  "section": null
}
\end{lstlisting}

\subsection{Difficultés d'extraction rencontrées}
\label{sec:bugs}

Trois familles de défauts, invisibles à la lecture rapide, ont été découvertes puis corrigées
(figure~\ref{fig:pipeline}). Elles méritent d'être détaillées car elles illustrent un point
méthodologique : \textbf{une vérification par échantillonnage ne suffit pas}. Aucun de ces trois
défauts n'a été trouvé à l'œil ; tous l'ont été par balayage systématique du corpus entier.

\begin{figure}[htbp]
  \centering
  % Chaîne d'extraction du corpus, et l'endroit exact où chacun des trois
% défauts de la section correspondante se manifeste.
\begin{tikzpicture}[font=\sffamily\scriptsize, node distance=6mm,
                    bloc/.append style={text width=23mm}]

\node[donnee, haut] (pdf) {PDF officiel\\472 pages};
\node[traitement, haut, right=of pdf] (txt)
  {Extraction du texte\\(PyMuPDF)};
\node[traitement, haut, right=of txt] (split)
  {Découpage par article\\expression régulière};
\node[traitement, haut, right=of split] (hier)
  {Reconstruction de\\la hiérarchie};
\node[donnee, haut, right=of hier] (out)
  {Corpus JSONL\\\NbArticles{} articles};

\draw[flux] (pdf) -- (txt);
\draw[flux] (txt) -- (split);
\draw[flux] (split) -- (hier);
\draw[flux] (hier) -- (out);

% ------------------------------------------- défauts, rattachés à leur étape
% Les pastilles se posent sur le coin du nœud, à cheval sur le bord, pour ne
% jamais recouvrir le texte qu'elles annotent.
\node[puce] at (txt.north east)   (p1) {1};
\node[puce] at (split.north east) (p2) {2};
\node[puce] at (hier.north east)  (p3) {3};

\node[annot, below=14mm of txt.south, text width=28mm, anchor=north]
  (a1) {\textbf{1.} L'exposant est perdu :\\
        «\,Article 334\textsuperscript{50}\,» ressort\\
        en \texttt{33450}.};
\node[annot, below=14mm of split.south, text width=28mm, anchor=north]
  (a2) {\textbf{2.} Un espace parasite\\
        casse l'ancrage : l'article\\
        135 disparaît en silence.};
\node[annot, below=14mm of hier.south, text width=28mm, anchor=north]
  (a3) {\textbf{3.} 32 titres se replient\\
        sur deux lignes : métadonnées\\
        tronquées, fin de titre\\
        absorbée par l'article.};

\draw[rejet, line width=0.5pt] (a1.north) -- (txt.south);
\draw[rejet, line width=0.5pt] (a2.north) -- (split.south);
\draw[rejet, line width=0.5pt] (a3.north) -- (hier.south);

% ------------------------------------------------------------ contrôle final
\node[bloc, draw=leafgreen, fill=leafgreen!8, text=leafgreen!45!black,
      text width=23mm, below=14mm of out.south, anchor=north] (ctrl)
  {Contrôle : séquence\\1--\NbArticles{} complète,\\empreinte SHA-256\\stable};
\draw[usage, draw=leafgreen!70] (out.south) -- (ctrl.north);

\end{tikzpicture}

  \caption[Chaîne d'extraction et défauts rencontrés]{Chaîne d'extraction du corpus, et l'étape à
  laquelle chacun des trois défauts se manifeste. Aucun ne produit d'erreur : ils dégradent le
  corpus en silence.}
  \label{fig:pipeline}
\end{figure}

\paragraph{Numéros d'articles corrompus par des appels de note}
PyMuPDF ne préserve pas la mise en exposant. Un appel de note collé à un numéro d'article dans le
PDF --- «~Article~334\textsuperscript{50}~» --- ressort donc comme une seule suite de chiffres sans
séparateur : \texttt{33450}. L'expression régulière, gourmande, consommait la suite entière et
produisait un article au numéro absurde. La correction consiste à plafonner la capture à trois
chiffres, ce qui est sûr puisque aucun article de cette loi ne dépasse \NbArticles{}.

\paragraph{Un article entièrement absent}
Le corpus comptait 588 articles au lieu de \NbArticles{}. Un espace parasite avant «~Article~135~»
dans le texte extrait cassait l'ancrage en début de ligne pour cette seule occurrence ; le contenu
de l'article était silencieusement absorbé dans le corps du précédent. Le défaut a été détecté en
cherchant les trous dans la séquence 1--\NbArticles{} --- un contrôle qui ne coûte rien et qu'il
aurait fallu écrire dès le premier jour.

\paragraph{Titres de hiérarchie repliés sur plusieurs lignes}
Trente-deux titres de section sont trop longs pour tenir sur une ligne du PDF et se replient sur
une deuxième, parfois une troisième. L'expression régulière ne capturant que la première ligne,
deux conséquences se cumulaient : les métadonnées de hiérarchie restaient tronquées pour
\emph{tous} les articles concernés, et la fin du titre fuyait dans le corps de l'article adjacent,
polluant le texte effectivement vectorisé.

\paragraph{La même cause, sur les en-têtes de hiérarchie}
Le plafonnement à trois chiffres corrige les \emph{numéros} d'articles, mais la perte de la mise
en exposant frappe aussi les titres. Un balayage systématique de tous les libellés de hiérarchie
du corpus --- et non des seuls articles --- a révélé quatre en-têtes portant encore un appel de
note soudé à leur dernier mot :

\begin{center}\small\ttfamily
\begin{tabular}{@{}l@{}}
Titre IV : De l'hygiène et de la sécurité des salariés\textcolor{anrtred}{37}\\
Section II : Du salaire minimum légal\textcolor{anrtred}{54}\\
Section III : De la saisie-arrêt\textcolor{anrtred}{60} et de la cession des salaires\\
Chapitre premier : Des agents chargés de l'inspection du travail\textcolor{anrtred}{88}\\
\end{tabular}
\end{center}

\noindent La correction retire tout groupe de un à trois chiffres soudé à une lettre dans un
en-tête. La règle est sûre : un titre qui numérote le fait avec une espace («~Section~2~») ou à
l'intérieur d'une référence de loi («~65-99~»), deux formes que le motif ne touche pas.

\begin{constat}
La régénération du corpus après cette correction ne modifie \textbf{aucun texte d'article} : elle
ne touche que 83 champs de métadonnées, sur les quatre libellés ci-dessus. Le contrôle est
important, car il établit que la correction ne pouvait pas déplacer les résultats du
chapitre~\ref{chap:evaluation} : ces champs ne sont pas vectorisés
(section~\ref{sec:indexation}), donc invisibles pour la recherche.

\smallskip
Ce défaut avait d'abord été documenté comme une limite assumée, au motif qu'y toucher obligerait à
rejouer toute l'évaluation. La vérification ci-dessus montre que cet argument était faux : le
champ n'étant pas vectorisé, aucun chiffre ne pouvait bouger. Le corriger ne coûtait donc rien.
\end{constat}

\subsection{Traitement d'un cas particulier : les articles abrogés}

Les articles~32 et~256 sont abrogés dans la version française de 2011 utilisée --- leurs
dispositions relatives au service militaire ayant été supprimées par la loi~48-06 --- mais ont été
rétablis par la loi~02.21 en 2021. Le texte de remplacement, vérifié contre la version arabe en
vigueur, a été réintégré dans le corpus.

Ces deux articles portent un indicateur explicite (\texttt{amende\_2021}) afin que cette
déviation par rapport au PDF source reste visible \emph{dans la donnée elle-même}, et pas seulement
dans la documentation. Le principe est général : une donnée qui s'écarte de sa source doit
transporter le signal de cet écart, faute de quoi il se perd à la première reprise du projet.

\subsection{Reproductibilité du corpus}

La chaîne d'extraction est déterministe : un clone vierge du dépôt suivi d'une réexécution du
parseur produit un corpus \textbf{identique octet pour octet} au corpus versionné. Cette propriété
a été vérifiée explicitement par comparaison d'empreintes SHA-256, et non simplement supposée.

\section{Indexation}
\label{sec:indexation}

Les \NbArticles{} articles sont vectorisés avec BGE-M3 puis stockés dans une collection Chroma
persistante. Seul le \emph{texte légal} est embarqué : la hiérarchie est conservée comme
métadonnée filtrable plutôt qu'incluse dans le vecteur.

Ce choix mérite une justification, car l'inverse est fréquent. Inclure «~Livre~II, Titre~III, De
la durée du travail~» en tête de chaque article ajoute à tous les articles d'une même section un
vocabulaire commun qui ne les distingue pas les uns des autres : le signal recherché par la
requête s'en trouve dilué, sans gain de précision.

Les vecteurs sont normalisés à l'indexation, ce qui permet de calculer la similarité cosinus par
simple produit scalaire, conformément à l'équation~\eqref{eq:cosinus}, indépendamment de la
métrique de distance configurée dans la base.

\section{Recherche hybride}

La recherche interroge les deux index puis fusionne les classements par RRF selon
l'équation~\eqref{eq:rrf}, avec une constante d'amortissement $k = 60$ et vingt candidats par
méthode.

\subsection{Une limite structurelle de la recherche par similarité}
\label{sec:refexplicite}

Un cas particulier a nécessité un traitement dédié. À la question «~Que dit l'article~32~?~», la
recherche sémantique échoue --- et elle échoue \emph{structurellement}, pas accidentellement. La
formulation de la question ne partage aucun vocabulaire avec le \emph{contenu} de l'article~32
(suspension du contrat, service militaire, grève) : la similarité entre les deux est donc faible,
alors même que l'article existe bel et bien dans le corpus. La recherche compare la question au
corps des articles, or une question \emph{sur} un article ne ressemble pas à cet article.

La correction consiste à détecter les références explicites («~article~N~») dans la question
elle-même et à garantir la présence de l'article correspondant dans le contexte, sans le faire
concourir dans un classement de similarité qu'il ne pouvait pas gagner. Le score d'abstention est
également forcé à 1,0 dans ce cas : un article nommé est par définition pertinent, et il serait
absurde que le garde-fou rejette une question dont on sait qu'elle porte sur un article réel du
corpus.

\begin{limite}
Ce chemin de code est resté \textbf{non testé pendant toute la durée du développement}. Le jeu de
référence d'origine ne comportait aucune question nommant un article par son numéro : la branche
existait, elle était raisonnée, mais rien ne prouvait qu'elle fonctionnait. Six questions ont été
ajoutées pour la couvrir (chapitre~\ref{chap:evaluation}).

\smallskip
Le forçage du score a par ailleurs une conséquence sur la mesure elle-même : les scores de
recherche de ces six questions valent 1,0 par construction et ne mesurent aucune similarité. Ils
sont donc exclus des statistiques de score du chapitre~\ref{chap:evaluation}, faute de quoi ils
gonfleraient artificiellement la bande «~droit du travail~».
\end{limite}

\section{Génération ancrée et garde-fous}

\subsection{Le prompt système}

La génération est encadrée par un prompt système imposant des règles strictes : répondre
uniquement à partir des articles fournis, citer systématiquement leurs numéros, employer une
phrase d'abstention exacte lorsque l'information est absente, et ne jamais donner de conseil
juridique personnalisé. Sa version finale est reproduite en annexe~\ref{ann:prompt}.

Trois règles n'y figuraient pas au départ. Elles ont été ajoutées en cours de projet, en réaction
à des défaillances observées lors de l'évaluation (section~\ref{sec:hallucination}) :
l'interdiction de compléter un aveu d'absence par une supposition, l'interdiction absolue de citer
un numéro de loi ou d'article provenant d'un autre code, et --- pour compenser l'effet des deux
premières --- une consigne de rééquilibrage énonçant que \emph{par défaut, le modèle répond}.

Le modèle est appelé à basse température (0,1), avec une pénalité de répétition et un plafond de
longueur. Ces deux derniers paramètres n'étaient pas prévus : ils ont été ajoutés après avoir
observé qu'une génération à basse température peut dégénérer en répétition indéfinie d'un même
mot, une fois le modèle à court de jetons dont il est confiant.

\subsection{Contrôle des citations}

Après génération, les numéros d'articles cités par le modèle sont extraits par expression
régulière, puis confrontés à la liste des articles réellement fournis. Tout numéro cité mais
absent du contexte est signalé.

Ce contrôle a lui-même fait l'objet d'une correction significative. Sa première version n'appariait
que la forme au singulier («~Article~152~») et manquait donc \emph{toutes} les citations au pluriel
(«~Articles~154 et~156~»), pourtant très fréquentes dans les réponses réellement produites. Le
garde-fou était ainsi partiellement inopérant, et les métriques d'évaluation sous-estimaient le
taux de citation --- un cas où l'instrument de mesure était lui-même défectueux, et où l'erreur
allait dans le sens de la sévérité, donc ne se signalait pas comme un problème.

\begin{constat}
Un garde-fou non testé est une hypothèse, pas une garantie. Les deux garde-fous du système ont
chacun révélé un défaut à l'usage : le contrôle des citations ignorait le pluriel, et le seuil
d'abstention s'est révélé mal calibré à sa première valeur. Aucun des deux ne se serait manifesté
sans un banc de mesure.
\end{constat}

\section{Interfaces}

Deux interfaces partagent exactement le même moteur.

\begin{itemize}
  \item Une \textbf{interface en ligne de commande}, qui constitue le livrable de référence. Elle
        expose la totalité de la traçabilité : score de recherche, articles retrouvés, citations
        non vérifiées. C'est celle qui a servi à toutes les mesures.
  \item Une \textbf{interface web conversationnelle} (React), destinée à la démonstration. Elle
        affiche les réponses au fil de leur génération, met en évidence les citations dans le
        texte, et permet de déplier chaque article source avec son chemin hiérarchique complet.
\end{itemize}

La surface technique reste volontairement du côté de la ligne de commande : l'interface web ne
montre que la réponse, ses sources et l'avertissement d'usage
(figure~\ref{fig:interface}). Un score de similarité affiché à un citoyen ne l'aiderait pas à
décider s'il peut faire confiance à la réponse ; il lui donnerait une apparence de rigueur sans le
moyen de l'interpréter.

% Captures de l'interface web, prises sur le système en fonctionnement.
%
% Produites par eval/capture_ui.js, qui pilote un navigateur sans affichage,
% pose réellement les questions au moteur local et enregistre l'écran. Ce ne
% sont donc pas des maquettes.
%
% Recadrées sur la zone de conversation — la barre latérale et l'espace vide
% sont retirés — et posées pleine largeur : à cette taille le texte de
% l'interface reste lisible à l'impression.
\newcommand{\captureui}[3]{%
  % #1 fichier, #2 titre, #3 commentaire
  \begingroup
  \setlength{\fboxsep}{0pt}\setlength{\fboxrule}{0.5pt}%
  \noindent\fcolorbox{rule}{white}{%
    \includegraphics[width=\dimexpr\linewidth-2\fboxsep-2\fboxrule\relax]{#1}}%
  \par\vspace{3pt}
  {\sffamily\scriptsize\raggedright
   \textbf{\color{navy}#2}\; \textcolor{slate}{#3}\par}
  \endgroup}


\begin{figure}[htbp]
  \centering
  \captureui{Figures/captures/01-reponse-citee-crop.png}{(a)~~Réponse citée}{%
    les numéros d'articles sont mis en évidence dans le texte ; les sources apparaissent en
    pastilles sous la réponse, suivies de l'avertissement d'usage.}

  \vspace{10pt}

  \captureui{Figures/captures/02-source-depliee-crop.png}{(b)~~Source dépliée}{%
    un clic sur une pastille ouvre l'article : son chemin hiérarchique complet, puis son texte
    légal, tels qu'ils figurent dans le corpus.}

  \caption[Interface web : réponse citée et source dépliée]{Interface web en fonctionnement. Ces
  vues ne sont pas des maquettes : elles sont produites par \module{eval/capture\_ui.js}, qui
  pilote un navigateur sans affichage, pose réellement les questions au moteur local et
  enregistre l'écran. Chaque réponse visible a donc été générée par Mistral~7B au moment de la
  capture.}
  \label{fig:interface}
\end{figure}

Le passage de (a) à (b) matérialise ce que le chapitre~\ref{chap:fondements} appelait la
vérifiabilité : le lecteur ne reçoit pas seulement un numéro d'article, il peut ouvrir le texte
cité et le lire, avec sa position exacte dans la loi. C'est la propriété que le RAG apporte et
qu'un modèle seul ne peut pas offrir.

La troisième vue (figure~\ref{fig:interface-abstention}) montre le comportement inverse : celui
où le système refuse.

\begin{figure}[htbp]
  \centering
  \captureui{Figures/captures/03-abstention-crop.png}{(c)~~Abstention}{%
    hors périmètre, le système rend la phrase d'abstention exacte, sans source ni contenu inventé.}
  \caption[Interface web : abstention]{Abstention affichée à l'utilisateur. Aucune source n'est
  proposée, et le score de recherche qui a déclenché le refus n'est pas montré : il relève de la
  traçabilité technique, exposée en ligne de commande, pas de ce que le citoyen a besoin de lire.}
  \label{fig:interface-abstention}
\end{figure}

\section*{Conclusion}
\addcontentsline{toc}{section}{Conclusion}

La réalisation confirme un point souvent sous-estimé : la qualité du corpus conditionne celle des
réponses, et ses défauts ne lèvent aucune alerte. Les trois familles de bugs d'extraction, toutes
invisibles à la lecture rapide, auraient dégradé silencieusement le système si elles n'avaient pas
été cherchées méthodiquement --- et l'une d'elles laisse encore un résidu, signalé plus haut. Le
chapitre suivant présente l'évaluation quantitative du système ainsi construit.
